% Shared definitions, adapted from the StaleyNET paper.

\usepackage{amsmath,amssymb,amsthm,mathtools}
\usepackage{graphicx}
\usepackage{booktabs}
\usepackage{array}
\usepackage[table]{xcolor}
\usepackage{tikz}
\usetikzlibrary{arrows.meta,positioning,calc,fit,shadows.blur,backgrounds}
\usepackage[most]{tcolorbox}
\usepackage{enumitem}
\usepackage{microtype}
\usepackage{xspace}

\definecolor{accent}{HTML}{3B4BA8}
\definecolor{definitionfill}{HTML}{EEF0FB}
\definecolor{codefill}{HTML}{F6F7F9}
\definecolor{rulefill}{HTML}{E7F5EF}
\definecolor{muted}{HTML}{5F6475}

% Syntax colors for pseudocode.
\definecolor{synkeyword}{HTML}{3B4BA8}
\definecolor{synstring}{HTML}{2E7D4F}

% Material tonal containers for figures: a fill and the dark text color of the same hue.
\definecolor{bluefill}{HTML}{D6E3FF}   \definecolor{blueink}{HTML}{001B3E}
\definecolor{violetfill}{HTML}{EADDFF} \definecolor{violetink}{HTML}{21005D}
\definecolor{greenfill}{HTML}{C3EFCB}  \definecolor{greenink}{HTML}{00210B}
\definecolor{amberfill}{HTML}{FFDDB3}  \definecolor{amberink}{HTML}{2A1700}
\definecolor{redfill}{HTML}{FFDAD6}    \definecolor{redink}{HTML}{410002}
\definecolor{grayfill}{HTML}{E9EAF0}   \definecolor{grayink}{HTML}{1B1C21}
\definecolor{edgeink}{HTML}{44474E}

% Definitions in shaded boxes.
\theoremstyle{definition}
\newtheorem{definition}{Definition}[section]
\tcolorboxenvironment{definition}{colback=definitionfill,colframe=definitionfill,boxrule=0pt,arc=2pt,left=6pt,right=6pt,top=4pt,bottom=4pt}

% The example line inside a definition.
\newcommand{\example}{\par\smallskip\noindent\textit{Example.}\ }

% A highlighted observation.
\newtcolorbox{observation}{colback=rulefill,colframe=rulefill,boxrule=0pt,arc=2pt,left=6pt,right=6pt,top=4pt,bottom=4pt,
  title={\sffamily\bfseries\footnotesize KEY OBSERVATION},coltitle=accent,colbacktitle=rulefill,
  attach title to upper,after title={\par\smallskip}}

% Pseudocode in the style of Cormen, Leiserson, Rivest and Stein: one
% statement per numbered line, procedure names in small capitals, and a
% wider gap after the last line of each group (end such a line with \groupend).
\newcounter{codeline}
\ExplSyntaxOn
\seq_new:N \l__proc_words_seq
\NewDocumentCommand{\proc}{m}{
  \textsf{\bfseries
    \seq_set_split:Nnn \l__proc_words_seq { - } { #1 }
    \seq_map_indexed_inline:Nn \l__proc_words_seq
      {
        \int_compare:nT { ##1 > 1 } { - }
        \tl_head:n { ##2 }
        \scalebox{0.8}{ \text_uppercase:n { \tl_tail:n { ##2 } } }
      }
  }
}
\ExplSyntaxOff

% A procedure in a shaded box. Arguments: the signature, the input in words, the output in
% words, and the link to the Rust lines that implement it. A procedure is one column wide.
\newcounter{procnum}
\newenvironment{procedure}[4]{%
  \setcounter{codeline}{0}\refstepcounter{procnum}%
  \begin{tcolorbox}[enhanced,colback=codefill,colframe=codefill,boxrule=0pt,arc=4pt,left=8pt,right=6pt,top=5pt,bottom=4pt]%
  \small\linespread{1.15}\selectfont\def\procsignature{#1}%
  \noindent{\sffamily\footnotesize\color{accent}Procedure \theprocnum}\hfill{\footnotesize #4}\par\vspace{2pt}%
  {\footnotesize\color{muted}\noindent\textit{Input:} #2\par\noindent\textit{Output:} #3\par}\vspace{5pt}%
  \begin{tabbing}%
  \hspace*{1.5em}\=\hspace{1.2em}\=\hspace{1.2em}\=\hspace{1.2em}\=\hspace{1.2em}\=\hspace{1.2em}\=\kill
  \hspace*{1.5em}\procsignature\ \{\\[1pt]}{%
  \\[1pt]\hspace*{1.5em}\}%
  \end{tabbing}\end{tcolorbox}}
\newcommand{\li}{\stepcounter{codeline}\llap{\scriptsize\color{muted}\thecodeline\hspace{0.6em}}\>\>}
\newcommand{\groupend}{\\[0.6\baselineskip]}
\newcommand{\kw}[1]{{\color{synkeyword}\textbf{#1}}}
\newcommand{\cmt}[1]{\hfill{\color{muted}\itshape\footnotesize // #1}}
\newcommand{\str}[1]{{\color{synstring}\texttt{"#1"}}}
\newcommand{\var}[1]{\ensuremath{\mathit{#1}}}
\newcommand{\st}[1]{\textsf{#1}}
\newcommand{\nil}{\textsc{nil}}
% A Rust type in a procedure.
\definecolor{syntype}{HTML}{7B3FA0}
\newcommand{\ty}[1]{{\color{syntype}\textsf{#1}}}

% Links to exact lines of the harness in the public repository.
\newcommand{\repo}{https://github.com/arpanpathak/anti_ai_slop}
\newcommand{\commit}{main}

\newcommand{\srcref}[3]{\href{\repo/blob/\commit/harness/src/#1\#L#2-L#3}{\texttt{#1}:#2--#3}}

% Rust listings with syntax colors.
\usepackage{listings}
\definecolor{synmacro}{HTML}{A1520B}
\lstdefinestyle{rust}{
  basicstyle=\fontsize{6.5}{7.8}\selectfont\ttfamily,
  keywords={fn,let,mut,match,loop,enum,impl,struct,self,Self,return,if,pub,const,use,mod},
  keywordstyle=\color{synkeyword}\bfseries,
  keywords=[2]{State,Result,Error,Output,Path,Command,OsStr,Answer,Harness,String,Ok,Err,Some},
  keywordstyle=[2]\color{syntype},
  keywords=[3]{bail,ensure,format,println,eprintln,json},
  keywordstyle=[3]\color{synmacro},
  comment=[l]{//},
  commentstyle=\color{muted}\itshape,
  morestring=[b]",
  stringstyle=\color{synstring},
  showstringspaces=false,
  columns=fullflexible,
  keepspaces=true,
  numbers=left,
  numberstyle=\tiny\color{muted},
  numbersep=5pt,
  xleftmargin=10pt,
  breaklines=true,
  breakindent=1.5em,
  postbreak=\mbox{{\color{muted}$\hookrightarrow$}\space}
}

% A one-column listing copied from a source file, with a link to the same lines on GitHub.
% Arguments: file in harness/src/, first line, last line, label, one-line title.
\newcounter{listnum}
\newcommand{\listingcard}[5]{%
  \begin{procfloat}[!t]%
  \begin{tcolorbox}[enhanced,colback=codefill,colframe=codefill,boxrule=0pt,arc=4pt,left=6pt,right=6pt,top=5pt,bottom=3pt]%
  \refstepcounter{listnum}\label{#4}%
  {\sffamily\footnotesize\color{accent}Listing \thelistnum}\hfill{\footnotesize\srcref{#1}{#2}{#3}}\par
  {\small #5}\par\vspace{2pt}%
  \lstinputlisting[style=rust,firstline=#2,lastline=#3,firstnumber=1]{../../harness/src/#1}%
  \end{tcolorbox}%
  \end{procfloat}}

% Procedures float in one column, like figures.
\makeatletter
\newcommand\fps@procfloat{tbp}
\newcommand\ftype@procfloat{4}
\newcommand\ext@procfloat{lop}
\newenvironment{procfloat}{\@float{procfloat}}{\end@float}
\makeatother

% Figures: Material cards with tonal fills, no outlines, a soft shadow below each box, and
% sharp arrowheads. The shadow is offset further down than it is blurred, so it never shows
% above a box as a ring.
\tikzset{
  figure/.style={font=\sffamily\small, >={Stealth[length=2.5mm, width=1.9mm]}, line width=0.75pt},
  elevated/.style={blur shadow={shadow xshift=0pt, shadow yshift=-1.6pt,
    shadow blur radius=1.4pt, shadow blur steps=14, shadow opacity=40,
    shadow blur extra rounding=1.4pt}},
  block/.style={rounded corners=6pt, draw=none, align=center, inner sep=5pt,
    minimum height=2.3em, elevated},
  chip/.style={rounded corners=3pt, draw=none, align=center, inner sep=2pt,
    font=\sffamily\scriptsize, minimum height=1.55em},
  region/.style={rounded corners=10pt, inner sep=8pt},
  region label/.style={font=\sffamily\scriptsize\bfseries, text=edgeink},
  edge label/.style={font=\sffamily\footnotesize, text=edgeink, align=center},
  flow/.style={->, draw=edgeink, line cap=butt},
}

% A figure card: a title and a one-line lede inside the figure, then the drawing.
\newcommand{\figcard}[3]{%
  \begin{tcolorbox}[enhanced,colback=white,colframe=white,boxrule=0pt,arc=6pt,left=7pt,right=7pt,top=6pt,bottom=7pt,
    drop fuzzy midday shadow=black!30]%
  {\sffamily\bfseries\small #1}\par\vspace{1pt}{\sffamily\footnotesize\color{muted}#2}\par\vspace{7pt}%
  \centering #3%
  \end{tcolorbox}}

\newcommand{\code}[1]{\texttt{#1}}

% A named constant in a procedure and in the text.
\newcommand{\const}[1]{\textsf{#1}}
